% Test LaTeX file for math_tree.py import-latex (F-MATH6)
% Uses Lean Blueprint \uses{} pattern for dependency declaration

\begin{definition}[Metric space]\label{def:metric}
A metric space is a set $X$ together with a function $d: X \times X \to \mathbb{R}$ satisfying
positivity, symmetry, and the triangle inequality.
\end{definition}

\begin{definition}[Open ball]\label{def:open-ball}
\uses{def:metric}
For $x \in X$ and $r > 0$, the open ball $B(x,r) = \{y \in X : d(x,y) < r\}$.
\end{definition}

\begin{definition}[Open set]\label{def:open-set}
\uses{def:open-ball}
A subset $U \subseteq X$ is open if for every $x \in U$ there exists $r > 0$ such that $B(x,r) \subseteq U$.
\end{definition}

\begin{definition}[Convergent sequence]\label{def:convergent}
\uses{def:metric}
A sequence $(x_n)$ converges to $x$ if for every $\epsilon > 0$ there exists $N$ such that $d(x_n, x) < \epsilon$ for all $n \geq N$.
\end{definition}

\begin{definition}[Cauchy sequence]\label{def:cauchy}
\uses{def:metric}
A sequence $(x_n)$ is Cauchy if for every $\epsilon > 0$ there exists $N$ such that $d(x_m, x_n) < \epsilon$ for all $m, n \geq N$.
\end{definition}

\begin{definition}[Complete metric space]\label{def:complete}
\uses{def:cauchy, def:convergent}
A metric space is complete if every Cauchy sequence converges.
\end{definition}

\begin{theorem}[Banach fixed-point theorem]\label{thm:banach}
\uses{def:complete, def:metric}
Let $(X, d)$ be a complete metric space and $T: X \to X$ a contraction mapping. Then $T$ has a unique fixed point.
\end{theorem}

\begin{corollary}[Picard iteration]\label{cor:picard}
\uses{thm:banach}
The sequence $x_{n+1} = T(x_n)$ converges to the unique fixed point of $T$.
\end{corollary}

\begin{definition}[Compact set]\label{def:compact}
\uses{def:open-set}
A subset $K \subseteq X$ is compact if every open cover has a finite subcover.
\end{definition}

\begin{theorem}[Heine-Borel theorem]\label{thm:heine-borel}
\uses{def:compact, def:metric}
In $\mathbb{R}^n$, a set is compact if and only if it is closed and bounded.
\end{theorem}

\begin{lemma}[Bolzano-Weierstrass]\label{lem:bolzano}
\uses{def:compact, def:convergent}
Every bounded sequence in $\mathbb{R}^n$ has a convergent subsequence.
\end{lemma}

\begin{theorem}[Extreme value theorem]\label{thm:evt}
\uses{def:compact, def:convergent}
A continuous function on a compact set attains its maximum and minimum.
\end{theorem}
